\documentclass[11pt,a4paper]{article}
\usepackage[a4paper,margin=2.5cm]{geometry}
\usepackage{fontspec}
\setmainfont{Latin Modern Roman}
\usepackage[hidelinks]{hyperref}
\usepackage{microtype}
\setlength{\parindent}{0pt}
\setlength{\parskip}{6pt}
\hypersetup{pdftitle={SE302 Requirements Document -- Team 2},pdfauthor={SE302 Section 2 Team 2}}
\newcommand{\reqitem}[3]{\par\vspace{2pt}\noindent\textbf{#1:} #2\par\nopagebreak
  \begin{list}{}{\setlength{\leftmargin}{1.2em}\setlength{\topsep}{0pt}\setlength{\itemsep}{0pt}\setlength{\parsep}{0pt}}\item\emph{Rationale:} #3\end{list}}
\begin{document}
\begin{center}
{\LARGE\bfseries Software Requirements Specification for\\[2pt] ``Exam Scheduler''}\\[10pt]
SE 302 -- Section 2, Team 2\\
Toprak Ege Güldoğan, Elif Eylül Tütüncü, Arda Korkmaz,\\
Arda Baran Günhan, Gökberk Erken\\
Milestone 1 -- November 2026
\end{center}
\section{Introduction}
``Exam Scheduler'' is going to be a desktop application that will be used by the student affairs office to prepare the exam schedule of an exam period. The user will provide the classrooms with their capacities and the courses with the students registered to them, define the length of the exam period, and the application will produce an exam schedule that respects a set of constraints. The produced schedule can then be viewed per classroom, per student, per course and per day. It is planned as a standalone application for a single user; it does not depend on any other application or service.

In this document a slot is one of the fixed time intervals in which an exam can be held on a given day, and the exam period is the set of days (each having the same number of slots) in which all exams have to be scheduled. The requirements below describe \emph{what} the application shall do, not how it will be designed or implemented.

\section{User Requirements}
The user requirements can be detailed as follows.

\reqitem{Functional Requirement 1}{The user shall be able to import a file that contains the classrooms and their capacities.}{The classroom data is available to the student affairs as a file, so it is not necessary to enter it manually. The sample file contains one classroom per line as the classroom name and its seat capacity. The number of classrooms and the capacities may differ from the sample data.}
\reqitem{Functional Requirement 2}{The user shall be able to import a file that contains the courses and the students registered to each course.}{The scheduler needs to know which students take which courses in order to avoid conflicts. The sample file contains one course code per course together with the IDs of its registered students. The number of courses and the number of students per course may differ from the sample data.}
\reqitem{Functional Requirement 3}{The user shall be able to import, optionally, a file that contains the list of all students.}{The sample data contains the list of all student IDs in a separate file. The customer stated that the classroom file and the course registration file are the minimum input, so this file is not required. When it is imported, the registered students of the courses can be checked against it.}
\reqitem{Functional Requirement 4}{The user shall be able to add, edit and remove classrooms and their capacities in the application.}{The data is expected to be created, edited or removed. Importing files is the main way of entering data, but the user shall also be able to correct the classroom data without changing the imported file.}
\reqitem{Functional Requirement 5}{The user shall be able to add, edit and remove courses in the application.}{The user shall be able to correct the course data without changing the imported file, for example to add a course that is missing from it.}
\reqitem{Functional Requirement 6}{The user shall be able to add, edit and remove the registrations of students to courses in the application.}{The registrations of students change often and are the main input of the conflict checking, so the user shall be able to correct them without changing the imported file.}
\reqitem{Functional Requirement 7}{The user shall be able to import a changed file again and run the scheduling again.}{The user may modify the data files and wants to re-import them and re-run the program. The previously imported data of that kind is replaced by the newly imported data.}
\reqitem{Functional Requirement 8}{The user shall be able to define the exam period, that is, the number of days and the number of slots per day.}{The number of days and the number of slots per day are fixed inputs of the scheduling. They determine whether a valid schedule exists, therefore the user has to be able to change them, for example to extend the exam period when no valid schedule exists.}
\reqitem{Functional Requirement 9}{The user shall be able to generate an exam schedule from the current data and exam period.}{Producing the exam schedule is the main purpose of the application. The schedule assigns every course to a day, a slot and a classroom.}
\reqitem{Functional Requirement 10}{The user shall be able to view, for each classroom, the exams that are held in that classroom.}{The user needs to see how each classroom is used during the exam period.}
\reqitem{Functional Requirement 11}{The user shall be able to view, for each student, the exam schedule of that student.}{The user needs to check the exams of an individual student, for example to answer a question of that student.}
\reqitem{Functional Requirement 12}{The user shall be able to view, for each course, the day and slot of its exam and the classroom assigned to it.}{This is the basic information that has to be announced for every exam.}
\reqitem{Functional Requirement 13}{The user shall be able to view, for each day, the exams that are held in each slot of that day.}{The user needs an overview of the whole day, for example to plan the supervision of the exams.}
\reqitem{Functional Requirement 14}{The user shall be informed when no valid schedule can be produced for the current data and exam period.}{It is possible that no schedule satisfies the mandatory constraints. In this case the application shall not produce an invalid schedule. The user can then extend the exam period (Functional Requirement 8) and generate the schedule again, which is one of the allowed ways of continuing.}
\reqitem{Functional Requirement 15}{The user shall be able to access a help menu inside the application.}{The user requires help menus within the application, because the application is going to be used without any other manual.}
\reqitem{Non-functional Requirement 1}{A user who has read only the help menu of the application shall be able to import the sample files, generate an exam schedule and open each of the four views (per classroom, per student, per course and per day) without any other assistance.}{No additional training is expected from the user. The help menu shall explain the file formats, the steps to produce a schedule and the available views. The requirement can be verified by giving the help menu and the sample files to a person who has not used the application and checking whether this task is completed.}
\section{System Requirements}
The system requirements can be detailed as follows.

\reqitem{Functional Requirement 16}{The system shall not assign two exams of the same student to two consecutive slots of the same day.}{This is a mandatory constraint of the customer: a student shall have at least one slot of break between two exams. Consecutive slots are understood as slots with consecutive numbers on the same day. This constraint cannot be relaxed under any circumstance.}
\reqitem{Functional Requirement 17}{The system shall not assign more than two exams of the same student to the same day.}{This is the second mandatory constraint of the customer. This constraint cannot be relaxed under any circumstance.}
\reqitem{Functional Requirement 18}{The system shall not assign two exams of the same student to the same slot.}{A student cannot attend two exams at the same time. Courses that have at least one common student therefore have to be scheduled to different slots.}
\reqitem{Functional Requirement 19}{The system shall not assign more than one exam to the same classroom in the same slot.}{A classroom can host only one exam at a time. The scheduler assumes that two different courses do not share a classroom in the same slot.}
\reqitem{Functional Requirement 20}{The system shall schedule every course exactly once, to one day, one slot and one classroom.}{Every course has one exam. In the sample data there are 20 courses, and each of them has exactly one exam.}
\reqitem{Functional Requirement 21}{The system shall assign every course to a classroom whose capacity is not smaller than the number of students registered to that course.}{All registered students of a course must have a seat. In the sample data each classroom has 40 seats and each course has 40 students.}
\reqitem{Functional Requirement 22}{The system shall search for a schedule that satisfies all the constraints above and shall not produce a schedule that violates any of them.}{The customer needs a valid solution; the best solution is not required. The mandatory constraints are never relaxed to obtain a solution. When the system cannot find a valid schedule, the user is informed as described in Functional Requirement 14.}
\reqitem{Functional Requirement 23}{The system should optionally improve the produced schedule according to a selected goal, such as spreading the exams equally over the exam period, finishing the exams in the minimum number of days, using the minimum or a balanced number of classrooms per day, or placing crowded, important or difficult courses at the beginning or the end of the exam period.}{The customer stated that optimizing the schedule is not a must and that it is up to the team. These goals are possible extensions on top of a valid schedule.}
\reqitem{Functional Requirement 24}{The system shall reject an imported file that is malformed or inconsistent and shall show a clear error message that describes the problem.}{The data comes from files prepared by the user. An incorrect file, for example a missing capacity, an unknown student or a duplicated course, shall not cause wrong schedules or crash the application.}
\reqitem{Non-functional Requirement 2}{The system shall be able to run on a Windows system, and it shall require only a Java Runtime Environment (JRE) to be installed.}{The target platform is a Windows desktop computer. The user cannot be asked to install any additional software, such as a database server.}
\end{document}
